%! Functions and Function Notation — Unit 1, Lesson 1.1 — Guided Notes (ANSWER KEY)
% THE in-class document, in the MAIN IDEAS / NOTES shape (LESSON_SHAPE.md §1): the vocabulary
% box, the hook, then ONE two-column guidednotes table. Each row is one idea — a label on the
% left; on the right one or two COMPLETE printed sentences, ONE large pre-drawn display, then a
% two-across grid of numbered problems with real writing room. Problems run 1..18.
% DENSITY: a \blank{} only where a word IS the answer; no mid-sentence blanks; two problems
% across; 2–2.6cm of answer space each.
%   I DO   : the teacher reads the definition, marks up the display, works the first problem
%            of each grid (1, 3, 7, 11).
%   WE DO  : the class works the rest of each grid, students holding the pen, every one
%            cold-called.
%   YOU DO : the last row — problems 15–18, alone, for 13 minutes while the teacher circulates.
% The crux is row 4, PIECEWISE: problem 13 asks for f(3) where the top rule REACHES x = 3 but
% does not OWN it — the open circle. Problem 16 transfers it to a formula with no graph.
%
% PAGE LOCKSTEP with notes_key/: the vocabulary write lines -> \ansline, \blank -> \ans, and
% the answer argument of every \pcell, \writespace and \labelbox filled. Nothing else differs.
% Inside a cell `\\` ENDS THE ROW — break lines with \par. (`\\` inside cases/tabular is safe:
% those environments scope it.) A row cannot break across pages: the figure and EACH grid row
% are separate sub-rows (`\\` then `& ...`, probgrid*).
\documentclass[10pt]{article}
\usepackage{precalculus-article}
\usepackage{precalculus-key}
\newcommand{\TallMath}[1]{$\displaystyle #1\rule[-1.4em]{0pt}{3.2em}$}

\begin{document}

\pageheader{Unit 1, Lesson 1.1}{Guided Notes --- Answer Key}

\vspace{0.12in}

\boxguard
\begin{vocabbox}
Fill in each term as we define it together.
\par\vspace{2pt}
\noindent\textbf{\textcolor{plum}{Function notation:}}\\[1pt]
\ansline{The symbol $f(x)$: $f$ is the rule's name, $x$ the input, $f(x)$ the output.}\\[3pt]
\noindent\textbf{\textcolor{plum}{Evaluate:}}\\[1pt]
\ansline{Find the output a function produces for a given input.}\\[3pt]
\noindent\textbf{\textcolor{plum}{Domain:}}\\[1pt]
\ansline{The set of all inputs a function uses.}\\[3pt]
\noindent\textbf{\textcolor{plum}{Range:}}\\[1pt]
\ansline{The set of all outputs a function produces; a repeat is listed once.}\\[3pt]
\noindent\textbf{\textcolor{plum}{Piecewise function:}}\\[1pt]
\ansline{One function built from several rules, each owning its own stretch of inputs.}\\[3pt]
\noindent\textbf{\textcolor{plum}{Undefined:}}\\[1pt]
\ansline{An expression with no value, such as division by zero --- that input is excluded.}
\end{vocabbox}

\vspace{0.1in}

\boxguard
\begin{hookbox}
A vending machine you can trust: press the same button, get the same item, every
time. Give the machine a name --- call it $P$ --- and let $P(b)$ be the price in
dollars of the item at button $b$. On the board:
\[ P(\text{B4}) = 1.25 \]

Say that equation out loud as an English sentence. What is the input here ---
and does an input have to be a number?
\ansline{``The item at button B4 costs \$1.25.'' The input is the button, B4 --- so no, an input does not have to be a number. It only has to be something the rule can take in and turn into exactly one output.}
\end{hookbox}

\small
\begin{guidednotes}
% ── ROW 1: function notation — naming the machine ────────────────────────────
\mainidea[Naming the machine]{Notation} &
A function has three parts, and the symbol $f(x)$ names all three at once: $f$ is the
\textbf{name} of the machine, $x$ is the \textbf{input} that goes in the slot, and $f(x)$ is
the \textbf{output} that comes out. It is read ``$f$ of $x$'' --- \textbf{never} ``$f$ times
$x$.'' The parentheses hold the input; they do not multiply.

\vspace{4pt}
A rideshare app charges a flat fee plus a rate per mile: $F(m) = 2.50 + 1.75m$.
\begin{center}
\begin{tikzpicture}[scale=0.95, >=stealth]
  \node[font=\small] (in) at (0,0) {input $m$};
  \node[draw=plum, very thick, fill=lilac, rounded corners=3pt, minimum width=4.2cm,
        minimum height=1.5cm, align=center, font=\small] (box) at (4.2,0)
        {\textbf{\Large\color{plum} $F$}\\ multiply by 1.75, then add 2.50};
  \node[font=\small] (out) at (8.6,0) {output $F(m)$};
  \draw[->, very thick, plum] (in) -- (box);
  \draw[->, very thick, plum] (box) -- (out);
  \node[font=\scriptsize\itshape, above] at (1.55,0.12) {goes in};
  \node[font=\scriptsize\itshape, above] at (7.05,0.12) {comes out};
\end{tikzpicture}
\end{center}
\vspace{2pt}
$m$ stands for \labelbox{2.8cm}{miles driven}\qquad $F(m)$ stands for \labelbox{3.2cm}{the fare, in \$}
\\
& \notesprompt{Read each statement as a sentence: name the input, the output, and the point on the graph.}
\begin{probgrid}{|Y|Y|}
\pcell{1}{$f(3) = 7$}{2.0cm}{Input 3, output 7; the point $(3, 7)$ is on the graph of $f$.} &
\pcell{2}{$F(4)$, for the rideshare fare above. What is it asking for? Find it.}{2.0cm}{The cost of a 4-mile ride. $F(4) = 2.50 + 1.75(4) = 9.50$; the point $(4, 9.50)$.} \\ \hline
\end{probgrid}
\\ \hline
% ── ROW 2: evaluating from a formula — the empty slot ────────────────────────
\mainidea[From a formula]{Evaluating} &
To \textbf{evaluate} a function, drop the input into the slot --- \emph{always inside
parentheses} --- and simplify. Evaluating never changes the rule; it just runs the machine once.
\textbf{Write the parentheses first, empty, then fill them.} That habit is what saves you when
the input is negative.

\vspace{4pt}
\begin{center}
\begin{tikzpicture}[scale=0.95, >=stealth]
  \node[font=\large] (rule) at (0,0) {$f(x) = 5x - 8$};
  \node[font=\large] (slot) at (5.6,0) {$f(\square) = 5(\square) - 8$};
  \draw[->, very thick, plum] (rule) -- node[above, font=\scriptsize\itshape] {empty the slot} (slot);
  \node[font=\large] (fill) at (5.6,-1.5) {$f(-2) = 5(-2) - 8$};
  \draw[->, very thick, plum] (slot) -- node[right, font=\scriptsize\itshape] {then fill it} (fill);
  \node[draw=goldacc, thick, fill=hookbg, rounded corners=2pt, font=\small, align=left]
        at (0,-1.5) {\textbf{The trap:} $(-4)^2$ and $-4^2$\\ are \emph{not} the same number.};
\end{tikzpicture}
\end{center}
\\
& \notesprompt{Evaluate. Show the slot filled before you simplify.}
\begin{probgrid}{|Y|Y|}
\pcell{3}{$f(x) = 5x - 8$. Find $f(3)$.}{2.0cm}{$f(3) = 5(3) - 8 = 15 - 8 = 7$} &
\pcell{4}{$f(x) = 5x - 8$. Find $f(-2)$.}{2.0cm}{$f(-2) = 5(-2) - 8 = -10 - 8 = -18$} \\ \hline
\end{probgrid}
\\
& \begin{probgrid}*{|Y|Y|}
\pcell{5}{$g(x) = x^2 + 3x$. Find $g(-4)$.}{2.0cm}{$g(-4) = (-4)^2 + 3(-4) = 16 - 12 = 4$} &
\pcell{6}{Compute $(-4)^2$ and $-4^2$. Which one did problem 5 need, and what made the difference?}{2.0cm}{$(-4)^2 = 16$ but $-4^2 = -16$. Problem 5 needed 16: the parentheses square the whole $-4$.} \\ \hline
\end{probgrid}
\\ \hline
% ── ROW 3: domain and range ──────────────────────────────────────────────────
\mainidea[Every input,]{Domain \& Range} &
The \textbf{domain} is the set of all inputs a function uses; the \textbf{range} is the set
of all outputs it produces, with a repeated output listed once. From a graph, read the domain
left to right and the range bottom to top; write either as an inequality, $0 \le x \le 8$, or
with brackets, $[\,0, 8\,]$ --- a square bracket means the endpoint \textbf{is} included.
\\
& \begin{center}
\begin{tikzpicture}[x=0.55cm, y=0.4cm, >=stealth]
  \draw[linegray, very thin, step=1] (0,0) grid (8,7);
  \draw[->, thick] (0,0) -- (8.6,0) node[below right] {\small $x$};
  \draw[->, thick] (0,0) -- (0,7.6) node[above] {\small $f(x)$};
  \foreach \x in {1,2,3,4,5,6,7,8} \node[below] at (\x,0) {\small \x};
  \foreach \y in {1,3,5,7} \node[left] at (0,\y) {\small \y};
  \draw[very thick, plum] (0,5) -- (2,1) -- (5,7) -- (8,4);
  \fill[plum] (0,5) circle (2.5pt);
  \fill[plum] (8,4) circle (2.5pt);
  % domain bracket under the x-axis
  \draw[very thick, goldacc, |-|] (0,-1.1) -- (8,-1.1);
  \node[below, font=\scriptsize\bfseries\color{goldacc}] at (4,-1.2) {left to right: every $x$ the graph uses};
  % range bracket left of the y-axis
  \draw[very thick, redacc, |-|] (-1.3,1) -- (-1.3,7);
  \node[rotate=90, above, font=\scriptsize\bfseries\color{redacc}] at (-1.5,4) {bottom to top: every height};
\end{tikzpicture}
\end{center}
\vspace{2pt}
Gold bracket: \labelbox{2.0cm}{domain}\qquad Red bracket: \labelbox{2.0cm}{range}
\\
& \notesprompt{State the domain and the range.}
\begin{probgrid}{|Y|Y|}
\pcell{7}{$f = \{(-2, 6),\ (0, 4),\ (3, 6),\ (5, 1)\}$. Domain and range, as lists.}{2.0cm}{Domain $\{-2, 0, 3, 5\}$; range $\{1, 4, 6\}$ --- the 6 is listed once.} &
\pcell{8}{The graph above: domain and range as inequalities, then again with brackets.}{2.0cm}{$0 \le x \le 8$ and $1 \le f(x) \le 7$; in brackets, $[\,0,8\,]$ and $[\,1,7\,]$.} \\ \hline
\end{probgrid}
\\
& \begin{probgrid}*{|Y|Y|}
\pcell{9}{$k(x) = \dfrac{8}{x + 3}$. Which input is forbidden, and why? Write the domain in words.}{2.2cm}{$x = -3$: it makes the denominator 0, which is undefined. The domain is every real number except $-3$.} &
\pcell{10}{In problem 7 the output 6 appears twice. How many times is it listed in the range? Is $f$ still a function?}{2.2cm}{Once. Yes --- two inputs may share an output; only one input with two outputs is forbidden.} \\ \hline
\end{probgrid}
\\ \hline
% ── ROW 4: THE CRUX — piecewise functions and the boundary ───────────────────
\mainidea[One function, several rules]{Piecewise} &
A \textbf{piecewise function} is one function built from two or more rules, each rule owning
its own stretch of inputs. The condition written beside a rule is its \textbf{address} --- it
says exactly which inputs that rule is responsible for, and no input belongs to two rules.

\vspace{3pt}
\stepnum{1}\ Find the stretch the input falls in. \quad \stepnum{2}\ Use \emph{only} that
rule. \quad \stepnum{3}\ Evaluate it the usual way. Never run both rules and pick the answer
you prefer --- the condition has already picked for you.
\\
& \begin{center}
\begin{tikzpicture}[x=0.55cm, y=0.4cm, >=stealth]
  \draw[linegray, very thin, step=1] (0,0) grid (7,7);
  \draw[->, thick] (0,0) -- (7.6,0) node[below right] {\small $x$};
  \draw[->, thick] (0,0) -- (0,7.6) node[above] {\small $f(x)$};
  \foreach \x in {1,2,3,4,5,6,7} \node[below] at (\x,0) {\small \x};
  \foreach \y in {2,4,6} \node[left] at (0,\y) {\small \y};
  % the boundary
  \draw[goldacc, very thick, dashed] (3,0) -- (3,7.2);
  \node[font=\scriptsize\bfseries\color{goldacc}, above right] at (2.35,7.2) {the boundary $x = 3$};
  % piece 1: 2x on [0,3)
  \draw[very thick, plum] (0,0) -- (3,6);
  \fill[plum] (0,0) circle (2.5pt);
  \draw[very thick, plum, fill=white] (3,6) circle (2.5pt);
  % piece 2: constant 2 on [3,7]
  \draw[very thick, plum] (3,2) -- (7,2);
  \fill[plum] (3,2) circle (2.5pt);
  \node[font=\scriptsize\itshape, right] at (3.25,6.6) {reaches $x=3$ but does not own it};
\end{tikzpicture}
\end{center}
\vspace{2pt}
\[ f(x) = \begin{cases} 2x & 0 \le x < 3 \\[2pt] 2 & x \ge 3 \end{cases} \]
\vspace{-4pt}

At $x = 3$ the rule that owns the input is the \labelbox{2.0cm}{bottom} one, so $f(3) = $ \labelbox{1.2cm}{$2$}
\\
& \fcolorbox{plum}{lilac}{\parbox{\dimexpr\linewidth-2\fboxsep-2\fboxrule\relax}{\centering
\textbf{An open circle marks a value the rule gets close to but never reaches. A filled dot
marks the one output the function actually has there. Every input gets exactly one output ---
so at a boundary, exactly one rule wins.}}}

\vspace{3pt}
\textbf{In your own words:} how do you decide which rule to use before you evaluate anything?
\writespace{1.4cm}{Read the conditions first and find which stretch the input falls in. That rule is the only one that applies --- evaluate it and stop.}
\\
& \notesprompt{Decide the rule first, then evaluate. Name the rule you used.}
\begin{probgrid}{|Y|Y|}
\pcell{11}{Using $f$ above, find $f(1)$. Which rule owns $x = 1$?}{2.0cm}{$x = 1$ lies in $0 \le x < 3$, so the top rule owns it: $f(1) = 2(1) = 2$.} &
\pcell{12}{Using $f$ above, find $f(5)$.}{2.0cm}{$x = 5$ satisfies $x \ge 3$, so the bottom rule: $f(5) = 2$.} \\ \hline
\end{probgrid}
\\
& \begin{probgrid}*{|Y|Y|}
\pcell{13}{Find $f(3)$. The top rule reaches $x = 3$ --- explain why 6 is still the wrong answer.}{2.2cm}{$f(3) = 2$. The top rule's address stops at $x < 3$, so 3 is not one of its inputs --- that is the open circle. $x = 3$ belongs to the bottom rule.} &
\pcell{14}{$g(x) = x^2$ when $x \le -1$, and $g(x) = 4 - x$ when $x > -1$. Find $g(-3)$ and $g(-1)$.}{2.2cm}{$-3 \le -1$, so $g(-3) = (-3)^2 = 9$. And $-1 \le -1$, so $g(-1) = (-1)^2 = 1$.} \\ \hline
\end{probgrid}
\\ \hline
% ── LAST ROW: YOU DO — alone, while the teacher circulates (13 min) ──────────
\mainidea[You do, alone]{Table \& Formula} &
A function $s$ given only by a table, a piecewise function $h$ given only by its rules, and
two functions given by a formula.
\\
& \begin{center}
\begin{tikzpicture}[scale=1.0]
  \node[font=\normalsize, anchor=north] at (0,1.5) {%
    \renewcommand{\arraystretch}{1.3}%
    \begin{tabular}{|c|c|c|c|c|c|}
    \hline $x$ & 1 & 2 & 3 & 4 & 5 \\ \hline $s(x)$ & 12 & 9 & 9 & 6 & 3 \\ \hline
    \end{tabular}};
  \node[draw=plum, very thick, fill=lilac, rounded corners=3pt, inner sep=6pt,
        font=\normalsize, anchor=north] at (0,-0.35) {%
    $h(x) = \begin{cases} 5 - x & x < 2 \\[2pt] 2x & x \ge 2 \end{cases}$};
\end{tikzpicture}
\end{center}
\\
& \notesprompt{On your own. Show your work.}
\begin{probgrid}{|Y|Y|}
\pcell{15}{$f(x) = 3x - 7$. Find $f(5)$ and $f(-1)$. Fill the slot first.}{2.2cm}{$f(5) = 3(5) - 7 = 8$; \ $f(-1) = 3(-1) - 7 = -10$.} &
\pcell{16}{Using $h$ above, find $h(-1)$, $h(2)$, and $h(6)$. Name the rule you used each time.}{2.2cm}{$h(-1) = 5 - (-1) = 6$ (top rule); $h(2) = 2(2) = 4$ (bottom, since $x \ge 2$); $h(6) = 2(6) = 12$ (bottom).} \\ \hline
\end{probgrid}
\\
& \begin{probgrid}*{|Y|Y|}
\pcell{17}{From the $s$ table: find $s(4)$, then state the domain and the range of $s$.}{2.2cm}{$s(4) = 6$. Domain $\{1, 2, 3, 4, 5\}$; range $\{3, 6, 9, 12\}$.} &
\pcell{18}{A garage charges \$3 per hour for the first 2 hours, then a flat \$6 for any longer stay, up to 12 hours. Write $C(h)$ as a piecewise function with both conditions.}{2.4cm}{$C(h) = 3h$ for $0 \le h \le 2$, and $C(h) = 6$ for $2 < h \le 12$.} \\ \hline
\end{probgrid}
\\ \hline
\end{guidednotes}

% ── THE NOTES END AT THE YOU DO ROW; AP Practice and the Homework follow ──────

\end{document}
